\chapter*{不定积分每日一题}

\refer{《积分的方法与技巧》金玉明\ 顾新身\ 毛瑞庭\ /\ 编著}

\exercise{求积分$\int{\frac{\d x}{a^2\sin^2{x}+b^2\cos^2{x}}}$}%在花括号内输入题目

\teacher{曹记铭}%在花括号内输入你的名字

\postdate{2022年12月19日}%在花括号内输入你想要发布那一天的日期

\result{\frac{1}{ab}\arctan{\qty(\frac{a}{b}\tan{x})+C}}%在花括号内输入本题结果

\analysis{
    先观察整体结构，我们发现分子为常数，分母是三角的二次.如果对这么一个庞然大物进行凑微，许多同学不知道如何下手，于是很多同学想到了万能公式，但次数更高后你会发现其实更麻烦了.如果联想到 是二次分式型，正好和分母形式相似，那么我们尝试对分母如下进行变形：
    \formula{
        \begin{aligned}
            \int{\frac{\d x}{a^2\sin^2{x}+b^2\cos^2{x}}}&=\int{\frac{\d x}{\cos^2{x}\qty(a^2\tan^2{x}+b^2)}}\\
            &=\int{\frac{\d\tan{x}}{a^2\tan^2{x}+b^2}}
        \end{aligned}
    }
    那么这样的话，我们就可以将$\tan{x}$看作整体，形式就比较清晰了
    
    令$t=\tan{x}$，则
    \formula{
        \begin{aligned}
            \int{\frac{\d\tan{x}}{a^2\tan^2{x}+b^2}}&=\int{\frac{\d t}{a^2t^2+b^2}}=\int{\frac{\d t}{b^2\qty(\dfrac{a^2}{b^2}t^2+1)}}\\
            &=\frac{1}{ab}\int{\frac{\d\qty(\dfrac{a}{b}t)}{\qty(\dfrac{a^2}{b^2}t^2+1)}}\\
            &=\frac{1}{ab}\arctan{\qty(\frac{a}{b}t)}+C\\
            &=\frac{1}{ab}\arctan{\frac{a}{b}\tan{x}}+C
        \end{aligned}
    }
    经验总结：分母三角（$\sin{x},\ \cos{x}$）的次数比分子的次数大二时，可以用尝试凑$\d\tan{x}$进行化简.
}%在花括号内输入分析